\chapter{Logging, Capture and Replay}\label{ch:ll:logging-capture-and-replay}

A regulator asks why one order left at 09:30:00.012345678 with that price and that size. There are two ways to answer. One is
to search the logs for the lines around that time and reconstruct, from what the programmers chose to print, what the strategy
probably saw. The other is to take the inputs the trading process received that morning, each with the time it received
it, feed them to the same binary on a desk, and stop at the order: the book, the parameters in force and every value the
decision used are there, exactly as they were, because the program is deterministic (chapter~20) and its inputs were kept.
This chapter builds the second answer and the machinery it needs: a log that costs a few nanoseconds on the hot path because
it writes raw numbers and formats them later, an \omterm{def:ll:logging-capture-and-replay:journal}{input journal} that records what arrived and when, and a replay driver that
turns a production day into a reproducible experiment.

\section{Why formatted logging is not for the hot path}

A formatted log line does a surprising amount of work: it parses a format string, converts each number to decimal text,
copies the result into a buffer, and, depending on the library, takes a lock and makes a system call. On the laptop, a line
with four numbers costs about \qty{100}{\nano\second} at the median through \texttt{snprintf} into a buffer, about as much
through \texttt{fprintf} into a fully buffered file, and \qty{210}{\nano\second} through a string stream
(\cref{fig:ll:logging-capture-and-replay:calls}); and its tail is long, because the file's buffer fills and is written, and
because the stream allocates. A strategy engine that handles an event in \qty{50}{\nano\second} (chapter~20) cannot afford one
such line per event, and a system whose logging is switched off in production to save that time has nothing to look at when
something goes wrong.

The answer is not to log less but to do less when logging. Everything in a log line is known when the program is compiled
except the numbers: the text, the positions of the arguments and their types. So the hot path can write only the numbers,
with an identifier for the statement, and leave the formatting to a program that runs later, elsewhere, when someone reads
the log. The \texttt{C++} standard library of the compiler used here (g++ 11) does not even provide \texttt{std::format}; it would not
change the conclusion, since formatting is the cost being removed.

\section{Binary logging and deferred formatting}

\begin{definition}[Binary log, deferred formatting]\label{def:ll:logging-capture-and-replay:binlog}
A \emph{binary log}\index{binary log} is a log whose records hold a statement identifier, a timestamp and the raw values of the
statement's arguments, not text; the statements' format strings are stored once, in the log's header or alongside the binary.
\emph{Deferred formatting}\index{deferred formatting} is the practice of turning such records into text only when the log is
read, by a separate program, so that the program that writes the log never converts a number to text.
\end{definition}

The idea is established: NanoLog, presented at the 2018 USENIX Annual Technical Conference, extracts the static parts of log
statements at compile time, writes a compacted \omterm{def:ll:logging-capture-and-replay:binlog}{binary log} at run time and re-inflates it offline, and reports an invocation
cost of 8 nanoseconds in \omterm{def:ll:cpu-microarchitecture:ubench}{microbenchmarks} and 18 in applications. The build's \texttt{firm.binlog} does the same with nothing
but the language. A statement is written \texttt{log.log<\textquotedbl{}fill \{\} of \{\} at \{\}, position \{\}\textquotedbl{}>(ts, id, qty, price, pos)}: the format
string and the argument types are template arguments, so the statement's identifier is computed at compile time, a hash of both
(\cref{lst:ll:logging-capture-and-replay:id}), and the statement registers its format in a table during static initialisation,
before \texttt{main} runs. At run time, a call writes the identifier, the timestamp and the arguments' bytes into a slot of a
single-producer ring (chapter~12), and returns (\cref{lst:ll:logging-capture-and-replay:log}). If the ring is full, the record
is dropped and counted: the hot path never waits for the log.

\omcode{code/firm/binlog/cpp/firm_binlog.hpp}{62}{70}{A statement's identifier, computed by the compiler from its format and
the types of its arguments.}\label{lst:ll:logging-capture-and-replay:id}

\omcode{code/firm/binlog/cpp/firm_binlog.hpp}{125}{141}{The hot path of a log call: encode the raw arguments behind a
16-byte header and write the record into the ring, or drop it.}\label{lst:ll:logging-capture-and-replay:log}

A writer thread on another core drains the ring into a file. The file starts with a header that lists every statement
(identifier, argument types, format), sorted by identifier, then the records as they came out of the ring; the Python decoder
reads the header, decodes each record with its statement's types and formats it. The format is simple enough to be written by
three implementations: the Python reference encoder writes the fixture's sample log, and the C++ logger, through its ring and
through its writer thread, and the Rust logger each write the same statements to the same bytes. A test also counts heap
allocations during logging and finds none, and fills a small ring on purpose to check that the overflow is dropped and counted.

\begin{omfigure}
\begin{tikzpicture}[font=\scriptsize,>=Stealth,box/.style={draw,rounded corners=2pt,minimum height=0.75cm,align=center,
  inner sep=3pt,fill=omDef!12}]
\node[box] (hot) at (0,0) {hot path\\ \texttt{log<\textquotedbl{}fill \{\} of \{\}\textquotedbl{}>}};
\node[box,fill=omMeth!15,minimum width=2.4cm] (ring) at (3.8,0) {ring (SPSC)\\ id $|$ len $|$ ts $|$ raw args};
\node[box] (wr) at (7.1,0) {writer thread\\ (another core)};
\node[box,fill=gray!10] (file) at (10.3,0) {file: statements\\ header $+$ records};
\node[box,fill=omThm!15] (dec) at (10.3,-1.9) {decoder (Python):\\ text, offline};
\node[box] (feed) at (0,-1.9) {inputs as received\\ (feed, reports, parameters)};
\node[box,fill=gray!10] (jr) at (3.5,-1.9) {input journal\\ recv ts $|$ event};
\node[box,fill=omThm!15] (rep) at (6.6,-1.9) {replay: same binary,\\ same inputs};
\draw[->,thick] (hot) -- node[above]{\qty{5}{\nano\second}} (ring);
\draw[->] (ring) -- (wr);
\draw[->] (wr) -- (file);
\draw[->] (file) -- (dec);
\draw[->] (feed) -- (jr);
\draw[->] (jr) -- (rep);
\draw[<->,dashed] (rep) -- node[below=2pt]{compare} (dec);
\end{tikzpicture}

\omcaption{Logging, capture and replay. The hot path writes raw records into a ring and returns; a writer thread drains them
to a file whose header holds every statement's format; formatting happens offline. The inputs are journaled as they arrive,
and a replay of the journal through the same binary produces the same log.}\label{fig:ll:logging-capture-and-replay:flow}
\end{omfigure}

\Cref{fig:ll:logging-capture-and-replay:calls} is the result: about \qty{5}{\nano\second} per call at the median, twenty
times less than \texttt{snprintf}, with a tail at the 99th percentile when the writer drains the ring and the \omterm{def:ll:memory-hierarchy-and-caches:line}{cache lines} it
reads move between the cores. The writer's manners matter. The first version polled the ring without pause, and every log
call paid for the \omterm{def:ll:memory-hierarchy-and-caches:line}{cache lines} the writer kept pulling away; the build's writer drains in batches and sleeps
\qty{100}{\micro\second} when there is little to drain (\cref{lst:ll:logging-capture-and-replay:writer}). The price is paid by
nobody on the hot path: the log reaches the file a fraction of a millisecond later.

\omcode{code/firm/binlog/cpp/firm_binlog.hpp}{206}{212}{The writer's loop: drain in batches, sleep when there is little to
drain.}\label{lst:ll:logging-capture-and-replay:writer}

\begin{omfigure}
\begin{tikzpicture}
\begin{axis}[width=10.5cm,height=5.6cm,ybar,bar width=12pt,ymode=log,ymin=1,ymax=2000,log origin=infty,
  xmin=-0.6,xmax=3.6,xtick={0,1,2,3},xticklabels={firm.binlog,snprintf,fprintf,ostringstream},
  ylabel={ns per call},tick label style={font=\small},label style={font=\small},grid=major,grid style={gray!25},
  legend pos=north west,legend style={font=\scriptsize},table/col sep=comma]
\addplot[fill=omThm!45,draw=omThm] table[x expr={floor(\coordindex/3)},y=ns,
  restrict expr to domain={mod(\coordindex,3)}{0:0}]{figdata/low-latency/23-logging-capture-and-replay/measured_logcall.csv};
\addlegendentry{median}
\addplot[fill=omMeth!40,draw=omMeth] table[x expr={floor(\coordindex/3)},y=ns,
  restrict expr to domain={mod(\coordindex,3)}{1:1}]{figdata/low-latency/23-logging-capture-and-replay/measured_logcall.csv};
\addlegendentry{99th percentile}
\end{axis}
\end{tikzpicture}

\omcaption{The cost of one log call carrying four numbers, by method (TSC-timed, 200\,000 calls each): the \omterm{def:ll:logging-capture-and-replay:binlog}{binary log} with
its writer thread draining the ring on another core, \texttt{snprintf} into a buffer, \texttt{fprintf} into a fully buffered
file, and a string stream. Measured on a laptop (Intel Core Ultra 7 155H) under WSL2. Data:
\texttt{bench\_log.py}.}\label{fig:ll:logging-capture-and-replay:calls}
\end{omfigure}

\section{Capture: inputs, outputs and time}

\begin{definition}[Input journal]\label{def:ll:logging-capture-and-replay:journal}
An \emph{input journal}\index{input journal} is the record, kept by a process, of every input it received, in the order it
received them, each with the time of its reception: market events, order reports, timer firings that depend on the clock,
parameter changes and control messages. Together with the program's version, it is enough to reproduce everything the
process computed.
\end{definition}

A log of decisions says what the program did; only a journal of inputs says why, because it lets the decision be made again.
The build's journal is deliberately dull: a short header, then fixed records of 56 bytes, the reception time followed by the
48-byte event of the \omterm{def:ll:the-feed-handler:handler}{feed handler} (chapter~18). It is written by the thread that receives the inputs, at the point where they
enter the engine, so that its order is the engine's order. Two things are easy to get wrong. The journal must hold what the
process received, not what the exchange sent: the exchange's recording has no gaps, no retransmissions and no local arrival
times, and a replay of it answers a different question. And the times must be the process's own, taken once, at reception,
with the clock of chapter~5, because they are inputs too: a strategy that ages its quotes reads them. The exchange's time
stamp and the receive time stamp of One Quant Book~1 (chapter~28) are both kept, one inside the event, the other in front of it.

\begin{dated}{2026-09}{How long and how precisely records are kept}\label{dat:ll:logging-capture-and-replay:records}
In the European Union (consulted September 2026), Commission Delegated Regulation (EU) 2017/589 requires firms engaged in
high-frequency algorithmic trading to record their orders and keep the records for five years from the submission of each
order (article 28). Commission Delegated Regulation (EU) 2017/574 requires members of a trading venue using a high-frequency
algorithmic trading technique to keep their business clocks within 100 microseconds of UTC, with a time stamp granularity of
1 microsecond or better, and to review their traceability to UTC at least once a year.
\end{dated}

Journals compress well, because consecutive events share most of their bytes: the same instrument, nearby prices and times,
sequence numbers that count by one. \Cref{tab:ll:logging-capture-and-replay:journal} measures it on the two streams of the
build: about 3.8 to 1 with the general-purpose compressor zlib at its default level, about 15 bytes per event.

\begin{table}[ht]
\centering\small
\begin{tabular}{lrrrrr}
\toprule
stream & events & recorded & journal & compressed & ratio \\
\midrule
large tick (the simulator's open) & 189\,875 & \qty{120}{\second} & 10.6 MB & 2.8 MB & 3.76 \\
small tick (synthetic) & 6\,000 & \qty{6}{\milli\second} & 0.34 MB & 0.08 MB & 3.96 \\
\bottomrule
\end{tabular}
\caption{The \omterm{def:ll:logging-capture-and-replay:journal}{input journal} of two event streams: 56 bytes per event raw, about 15 compressed with zlib at level 6. Data:
\texttt{fig\_journal.py} (deterministic).}\label{tab:ll:logging-capture-and-replay:journal}
\end{table}

\section{Deterministic replay of a production day}

\begin{definition}[Deterministic replay]\label{def:ll:logging-capture-and-replay:replay}
A \emph{deterministic replay}\index{deterministic replay} runs a program on an \omterm{def:ll:logging-capture-and-replay:journal}{input journal}, in the journal's order, with time
taken from the journal, and produces bit for bit the outputs that the program produced when the journal was recorded; it
requires the program to be deterministic (chapter~20) and the journal to be complete.
\end{definition}

The replay driver reads the journal, rebuilds the events and gives them to the strategy engine with the simulated clock: the
build's test journals the small-tick day, replays it, and obtains the committed output hash of chapter~20. A replay is also
fast, because nothing waits for the market: on the laptop, the engine replays the simulator's two busy minutes around the open,
189\,875 events, journal reading included, in under \qty{50}{\milli\second}, about 2\,800 times faster than they happened. A day
replays in seconds, which changes what can be done with it: every change to the code can be replayed against last week, and
every incident can be replayed as often as needed with a debugger attached.

When two runs that should agree do not, the logs find where. Each run logs its actions; the two logs are compared record by
record (\cref{lst:ll:logging-capture-and-replay:diff}), and the first record that differs points at the first event whose
handling differed. The build's test plants a difference, a parameter change at event 3\,000 that the second run receives and
the first does not, as a value read from outside the journal would; the comparison finds the first differing action,
number 540 of 1\,084, just after the planted event, and nothing before it. This is the procedure of chapter~20's weekend problem,
mechanised: the log says where to look, and the replay lets one look as often as needed.

\omcode{code/firm/binlog/firm_binlog.py}{99}{107}{Offline: format the records, and find the first record at which two logs
differ.}\label{lst:ll:logging-capture-and-replay:diff}

\section{What must be kept, and for how long}

Three things are kept. The \omterm{def:ll:logging-capture-and-replay:journal}{input journal}, because it is the only record from which everything else can be recomputed; the
decision log, because it answers most questions without a replay and checks the replay when one is made; and the program:
the exact binary, its configuration and the parameter snapshots, since a journal replayed through a different build answers
nothing. The regulator's clocks and retention periods (the dated box) set the minimum: time stamps traceable to UTC with the
granularity required, and five years of order records. The storage is modest by the standards of the data a quantitative firm
keeps anyway, as the weekend problem computes, and the cost of not keeping it is the first sentence of this chapter.

\section{Tutorial: log, journal, replay, bisect}\label{tut:ll:logging-capture-and-replay:tut}

\begin{tutorial}
\textbf{Goal.} Log from the hot path at a few nanoseconds per call, journal a day's inputs, replay them to the bit, and find a
planted difference between two runs. \textbf{End state:} \cref{fig:ll:logging-capture-and-replay:calls}, the journal table,
the replay's speed, and green tests in three languages.
\begin{tutsteps}
\item \textbf{Log.} The C++ test logs the fixture's statements through the ring and through the writer thread and compares the
bytes with \texttt{data/sample.blog}, which the Python reference wrote; the Rust test writes the same bytes; the Python decoder's \texttt{read} and \texttt{text} turn them into
\texttt{data/sample.txt}.
\item \textbf{Measure.} \texttt{python bench\_log.py} times a log call by each method and a replay of the open's journal.
\item \textbf{Journal and replay.} The test journals the small-tick day, replays it through \texttt{firm.stratengine}, and
checks the committed hash.
\item \textbf{Bisect.} \texttt{ll\_bisect\_test.cpp} runs the engine twice, plants a change in the second run, logs both, and
finds the first differing record.
\end{tutsteps}
\textbf{What to change next.} Replace the planted parameter change by a wall-clock read (chapter~20's defect) and watch the
first difference move from run to run; compress the journal as it is written, in the writer thread.
\end{tutorial}

\section{Build: the binary log and the journal}\label{bld:ll:logging-capture-and-replay:binlog}

\begin{build}
\bfield{Purpose} Every component of the trading path logs through it: the \omterm{def:ll:the-feed-handler:handler}{feed handler}'s gaps (chapter~18), the engine's
actions (chapter~20), the gateway's messages (chapter~21), the \omterm{def:ll:pre-trade-risk-and-kill-switches:gate}{risk gate}'s decisions (chapter~22); its journal is what the
resilient architecture of chapter~24 replicates and what the tests of chapter~25 replay.
\bfield{Interface} C++20 \texttt{firm::binlog}: \texttt{Logger(capacity)} with \texttt{log<\textquotedbl{}format\textquotedbl{}>(ts, arguments)},
\texttt{drain}, \texttt{dropped}; \texttt{Writer(logger, path)}; \texttt{header()}; \texttt{Journal} with \texttt{add} and
\texttt{records}. Python \texttt{firm\_binlog}: \texttt{site\_id}, \texttt{Writer}, \texttt{read}, \texttt{text},
\texttt{first\_difference}, \texttt{journal\_write}, \texttt{journal\_read}. Rust \texttt{firm\_binlog}: \texttt{Logger},
\texttt{read}, \texttt{text}, \texttt{site\_id}.
\bfield{Rules} No formatting, lock, system call or allocation on the hot path; a full ring drops and counts; statement
identifiers from the format and the types; the journal written where the inputs enter the engine, with reception times; the
same bytes from the three implementations.
\bfield{Acceptance tests} \texttt{code/firm/binlog/}: the fixture written byte for byte by Python, C++ (ring and writer thread)
and Rust, and decoded to its text; random statements round-tripped; identifiers that depend on the types, and collisions
refused; zero allocations while logging, and overflow counted; a journal replayed through the strategy engine to the committed
hash; the first difference between two logs found.
\bfield{Stretch} Compression in the writer thread; a statement per component with a severity filtered at compile time;
time stamps from the TSC with calibration records in the log, converted offline; a journal of the gateway's reports next to
the market data.
\end{build}

\begin{omsources}
\item S.~Yang, S.~J.~Park and J.~Ousterhout, ``NanoLog: a nanosecond scale logging system'', \emph{USENIX Annual Technical
Conference}, 2018.
\item Commission Delegated Regulation (EU) 2017/589 (RTS 6), article 28.
\item Commission Delegated Regulation (EU) 2017/574 (RTS 25), article 4 and annex.
\end{omsources}

\section{Exercises}

\begin{exercise}[$\star$]\label{exo:ll:logging-capture-and-replay:1}
How many bytes does a record of the statement \texttt{\textquotedbl{}fill \{\} of \{\} at \{\}, position \{\}\textquotedbl{}} with four 64-bit arguments take,
and how many does its formatted line take for the values 40, 100, 999\,900 and $-200$ with a 14-digit time stamp? What, then,
does the \omterm{def:ll:logging-capture-and-replay:binlog}{binary log} save?
\end{exercise}

\begin{exercise}[$\star$]\label{exo:ll:logging-capture-and-replay:2}
Why does the statement's identifier depend on the types of its arguments and not only on its format?
\end{exercise}

\begin{exercise}[$\star$]\label{exo:ll:logging-capture-and-replay:3}
The open's journal has 189\,875 records of 56 bytes plus a 12-byte header. Check its size, and the bytes per event after
compression, from \cref{tab:ll:logging-capture-and-replay:journal}.
\end{exercise}

\begin{exercise}[$\star\star$]\label{exo:ll:logging-capture-and-replay:4}
A component logs 2 million records a second during a burst; the writer thread is blocked for \qty{20}{\milli\second} by a slow
disk. How many ring slots are needed so that nothing is dropped, and how much memory is that with 64-byte slots?
\end{exercise}

\begin{exercise}[$\star\star$]\label{exo:ll:logging-capture-and-replay:5}
The hot path stamps records with the TSC rather than with \texttt{clock\_gettime}. What must the log also contain so that the
time stamps can be converted to UTC offline with a granularity of a microsecond?
\end{exercise}

\begin{exercise}[$\star\star$]\label{exo:ll:logging-capture-and-replay:6}
An engine handles an event in \qty{50}{\nano\second} and logs five lines per event. What does logging cost per event with the
\omterm{def:ll:logging-capture-and-replay:binlog}{binary log} and with \texttt{snprintf}, at the medians of \cref{fig:ll:logging-capture-and-replay:calls}?
\end{exercise}

\begin{exercise}[$\star\star\star$]\label{exo:ll:logging-capture-and-replay:7}
\emph{Coding.} Log the \omterm{def:ll:pre-trade-risk-and-kill-switches:gate}{risk gate}'s decisions (chapter~22) with a statement carrying the order's identifier, the code and the
mask; replay the gate's fixture, and decode the log into a count of refusals by reason that matches \texttt{expected.txt}.
\end{exercise}

\begin{exercise}[$\star\star\star$]\label{exo:ll:logging-capture-and-replay:8}
\emph{Find the flaw.} ``We do not need an \omterm{def:ll:logging-capture-and-replay:journal}{input journal}. We log every decision with its reason, and to replay a day we run the
strategy on the exchange's own recording of the market data.''
\end{exercise}

\section{Problem: To the Nanosecond}

\begin{problem}[{Weekend problem --- what it costs to be able to answer}]\label{pb:ll:logging-capture-and-replay:1}
A firm journals the inputs of its trading processes and must be able to replay any order of the last five years. Use
\cref{tab:ll:logging-capture-and-replay:journal} and the measured replay speed.

\medskip\noindent\textbf{Part I --- The rate.}
\begin{enumerate}
\item What is the average event rate of the open's stream?
\item At that average rate, how many events does a trading day of 6.5 hours hold?
\item How many bytes is that journal raw, and compressed at the measured ratio?
\item Why is an average over a busy open a pessimistic estimate for a whole day, and why is it still the right one to plan
with?
\end{enumerate}

\medskip\noindent\textbf{Part II --- Five years.}
\begin{enumerate}[resume]
\item With 252 trading days a year, how much does five years of such journals take, compressed?
\item The firm trades 200 instruments with similar streams. Scale the answer.
\item What else must be kept for a journal to be replayable five years later?
\item Why is the decision log, although smaller, not enough?
\end{enumerate}

\medskip\noindent\textbf{Part III --- Replay.}
\begin{enumerate}[resume]
\item At the measured speed, how long does a replay of one instrument's day take?
\item And of the 200 instruments, one after the other, on one core?
\item What limits the replay speed, and what does not?
\item Why must the replay read time from the journal rather than from the machine?
\end{enumerate}

\medskip\noindent\textbf{Part IV --- The verdict.}
\begin{enumerate}[resume]
\item State the \emph{named result}: journal storage per day at the measured rate and compression ratio, and replay speed
against real time.
\item What does a log call cost with \omterm{def:ll:logging-capture-and-replay:binlog}{deferred formatting}, and with \texttt{snprintf}?
\item What must the replay reproduce to answer the regulator's question, and how do you check that it did?
\item What happens to replayability when a component reads the wall clock?
\item Where in the process must the journal be written, and why there?
\item What does the planted-difference test demonstrate, and what would it take to find a real one?
\item What time stamp precision does the dated box require, and does a TSC-stamped log meet it?
\item In one sentence: what makes an order explainable years later?
\end{enumerate}
\end{problem}

\section{Interview questions}

\begin{interviewq}[$\star$ \iqroles{developer}]\label{iq:ll:logging-capture-and-replay:1}
Why is \texttt{printf}-style logging a problem on a hot path, and what do you do instead?
\end{interviewq}

\begin{interviewq}[$\star\star$ \iqroles{developer}]\label{iq:ll:logging-capture-and-replay:2}
Design a logger whose call costs under ten nanoseconds. What happens when its buffer is full?
\end{interviewq}

\begin{interviewq}[$\star\star$ \iqroles{developer}]\label{iq:ll:logging-capture-and-replay:3}
What would you record to be able to replay a trading day exactly, and where in the process would you record it?
\end{interviewq}

\begin{interviewq}[$\star\star$ \iqroles{developer}]\label{iq:ll:logging-capture-and-replay:4}
Two replays of the same day disagree. How do you find where?
\end{interviewq}

\begin{interviewq}[$\star\star$ \iqroles{developer, risk}]\label{iq:ll:logging-capture-and-replay:5}
A regulator asks why an order was sent three years ago. What do you need to answer, and how long does it take?
\end{interviewq}

\begin{interviewq}[$\star\star\star$ \iqroles{developer}]\label{iq:ll:logging-capture-and-replay:6}
Design capture and replay for a firm with dozens of trading processes on several venues. What is journaled, where, in what
format, for how long, and how is replay tested?
\end{interviewq}
